\documentclass{article}
\usepackage[utf8]{inputenc}
\usepackage[T1]{fontenc}
\usepackage{graphicx}
\usepackage{float}
\begin{document}

\section*{Slide 1}

Radiation   Studies   on   the   ATLAS   Link   and   Control 
Daughterboard   for   the   HL-LHC   Upgrade

1 st   Eduardo   Valdes   Santurio 
Fysikum,   Stockholm   University 
Stockholm,   Sweden 
eduardo.valdes@fysik.su.se 
ORCID:   0000-0001-9931-2896

2 nd   Samuel   Silverstein 
Fysikum,   Stockholm   University 
Stockholm,   Sweden 
silver@fysik.su.se 
ORCID:   0000-0001-7734-7617

3 rd   Christophe   Clement 
Fysikum,   Stockholm   University 
Stockholm,   Sweden 
christophe.clement@fysik.su.se 
ORCID:   0000-0003-3122-3605

4 th   Christian   Bohm 
Fysikum,   Stockholm   University 
Stockholm,   Sweden 
bohm@fysik.su.se

5 th   Holger   Motzkau 
Fysikum,   Stockholm   University 
Stockholm,   Sweden 
holger.motzkau@fysik.su.se 
ORCID:   0000-0002-0652-6236

multiplier   tubes   (PMTs)   to   the   off-detector   systems.   The   on- 
and   off   detector   electronics   are   interconnected   via   a   radiation 
tolerant   readout   link   and   control   Daughterboard   (DB).   Each 
DB is interfaced to a Upgrade TileCal Module that hosts up to 
12 PMT channels in the front-end by means of a 400-pin FMC 
connector.   A   total   of   896   DB   will   transmit   approximately   35 
Tbps of data via 5376 optical links to TileCal off-detector Pre- 
Processor   modules.   The   on-detector   systems,   which   include 
the DB, has been carefully engineered with redundancy layers 
and   radiation   mitigation   strategies   to   maintain   performance 
and   reliability   with   negligible   impact   on   data   taking   during 
LHC   runs.

Abstract ---We   present   a   summary   of   the   radiation   studies 
conducted   on   the   latest   version   of   the   ATLAS   Tile   Calorimeter 
(TileCal) readout and control Daughterboard (DB), developed for 
the   High   Luminosity   upgrade   of   the   Large   Hadron   Collider.   The 
DB functions   as a interface   hub   between   the   on-detector   and off- 
detector   electronics   by   to   handling   configuration   commands   and 
synchronous   timing   signals,   while   simultaneously   transmitting 
monitoring,   calibration,   and   physics   data   from   the   front-end 
electronics.   The   design   incorporates   Commercial-off-the-shelf 
(COTS)   devices   with   field-programmable   gate   arrays   (FPGAs) 
and   CERN   custom   radiation-hard   GBTx   ASICs.   The   COTS 
were   thoughtfully   studied   and   selected   for   their   performance 
in   radiation   environments   comparable   to   those   expected   over 
the   operational   lifetime   of   the   collider.   Each   DB   will   provide 
continous   readout   via   two   pairs   of   redundant   9.6   Gbps   uplinks 
continuously   receive   LHC   syncronous   timing   through   a   pair   of 
redundant 4.6 Gbps downlinks. The upgraded calorimeter system 
will   include   896   DBs   serving   a   total   of   5,376   high-speed   optical 
links.   A   double-redundant   architecture   has   been   implemented   to 
minimize single points of failure, while data integrity is preserved 
through   techniques   such   as   electronic   redundancy   layers,   Triple 
Mode   Redundancy   (TMR),   Soft   Error   Management   (SEM)   and 
the   GBT   protocol   with   Forward   Error   Correction   (GBT-FEC). 
Extensive   testing   was   performed   to   assess   the   impact   of   total 
ionizing   dose   (TID),   non-ionizing   energy   losses   (NIEL),   and 
single-event   effects   (SEE)   on   system   performance   under   realistic 
operating   conditions. 
Index Terms ---Radiation Tolerant, TID, NIEL, SEU, SEL, SEE

II.   T HE  D AUGHTERBOARD   DESIGN

The   DB   is   a   read-out   link   and   control   board,   designed   for 
the   upgrade   of   TileCal   [2].   A   set   of   commercial   of   the   shelf 
components   and   the   CERN   GBTx   custom   rad-hard   ASIC   [4] 
was   selected   for   the   design   to   achieve   good   performance   for 
a   decade   of   full   functionality   in   the   relatively   low   radiation 
environment   simulated   at   the   geographical   position   of   the   DB 
in   the   ATLAS   detector.   The   design   follows-up   a   redundancy 
layer   of   two   functionally-equivalent   halves   used   on-detector 
in   TileCal   modules   [1]. 
Each   DB   powers   four   SFP+   modules   that   drive   two   4.8 
Gbps   downlinks   and   four   9.6   Gbps   uplinks.   A   pair   of   CERN 
GBTx   hard   ASICs   receive   LHC   synchronized   clock   signals 
and   configuration   commands   through   the   downlinks.   Mean- 
while,   a   pair   of   Kintex   Ultrascale   (KU)   FPGAs   distribute 
clocks   and   configuration   received   from   the   GBTx   to   the 
front-end   and   transmits   front-end   data   and   monitoring   signals 
through   the   uplinks.   The   ASICs,   flash   memories   and   FPGAs 
of   the   whole   board   can   be   remotely   configured   through   the 
downlinks. 
The   xADCs   of   the   KU   FPGAs   monitor   temperatures,   volt- 
age   stability   and   current   from   both   the   DB   and   the   front-end 
to   assure   that   the   on-detector   electronics,   inaccessible   during 
runs, is well protected against any over-current or Single Event

I.   I NTRODUCTION

The   ATLAS   Tile   Calorimeter   (TileCal)   is   in   charge   of 
identifying   hadronic   jets   and   measuring   transverse   missing 
energy   [1].   TileCal   will   be   impacted   by   a   harsher   radiation 
environment   and   increased   rates   of   pileup   as   a   consequence 
of   the   High   Luminosity   Large   Hadron   Collider   (HL-LHC) 
programme.   An   upgrade   of   the   legacy   calorimeter   electronics 
is   needed   to   guarantee   full   functionality   over   the   10   year 
period of the HL-LHC. The upgraded electronics will continu- 
ously transmit digitized data from approximately 10000 photo-

Work   supported   by   Stockholm   University   and   CERN.

\section*{Slide 2}

some of the samples being able to reach 21  $\times$  10 12   n/cm 2   and 
more   than   200   Gy   without   any   failures. 
For   SEEs   three   different   tests   were   performed   for   studying 
the   SELs   and   Single   Event   Upsets   (SEUs).   A   KU+   FPGA 
was   irradiated   to   2.2   $\times$   10 11   p/cm 2   with   54   MeV   protons, 
where   a   SEU   rate   of   24   SEU/( 10 9   p/cm 2 )   and   an   SEL   rate   of 
2   SEL/( 10 9   p/cm 2 ).   The   presence   of   SEL   in   the   KU+   FPGA 
disqualified it for usage on any on-detector design. In the case 
of   the   KU   FPGA   two   tests   were   done   with   54   MeV   protons 
and   226   MeV   protons   where   no   SEL   wee   observed   up   to 
2.62   $\times$   10 11   p/cm 2   and   3.2   $\times$   10 11   p/cm 2   respectively.   As 
for   the   SEU   rates,   we   measured   116   SEU/( 10 9   p/cm 2 )   and   80 
SEU/( 10 9   p/cm 2 ). Worst case scenario, the SEU rates represent 
roughly 5498 SEUs per year out of which approximately 5487 
are correctable by the SEM and wont impact the nominal runs 
due to the use of TMR. The remaining 11 uncorrectable errors 
per year will be managed by the TMR and can be removed by 
reloading   the   FPGA   configuration   from   the   FLASH   memory 
during   data   acquisition   stops.

Latch-ups   (SEL).   Additionally,   headers   were   added   for   extra 
flexibility in digital and analog monitoring capabilities. All the 
components   utilized   in   the   design   were   selected   and   tested   to 
the   simulated   radiation   levels   for   the   HL-LHC   lifetime   plus 
safety   factors.   The   single   point   of   failure   mitigation   strategies 
to   reduce   the   impact   radiation   effects   include:   redundancy 
layers,   GBTx   ASICs,   Triple   Mode   Redundancy   (TMR),   the 
Xilinx   Soft   Error   Management   (SEM),   and   GBT-FEC   proto- 
col.

III.   T HE   SIMULATED   RADIATION   ENVIRONMENT

The 
ATLAS 
Radiation 
Tolerant 
Electronics 
Proposed 
Guideline   on   COTS   Lot   to   Lot   Variation   [4]   was   put   together 
by   the   collaboration   as   a   policy   for   better   consistency   in   the 
radiation   qualification   tasks   within   the   different   sub-detector 
electronics   designs   [4].   Due   to   the   cylindrical   symmetry   of 
the ATLAS detector the radiation requirements were extracted 
from   the   radiation   levels   of   the   DBs   positioned   in   the   most 
heavily   irradiated   environments   [5].   The   minimum   Total   Ion- 
izing   Dose   required   for   qualification   of   each   of   the   design 
components   id   108   Gy   of,   with   inclusion   of   safety   factors   of 
1.5   for   simulation   and   3.0   for   lot   variation   in   the   components. 
The   Non-Ionizing   Energy   Losses   (NIEL)   fluence   required 
for   qualification   is   13.16   $\times$   10 12   n/cm 2 ,   taking   into   account 
safety   factors   of   1.5   for   simulation,   1.3   when   using   a   non- 
monoenergetic   neutron   source,   and   3.0   to   account   for   lot 
variation   of   the   components.   The   total   expected   fluence   for 
Single   Event   Effects   (SEE)   is   4.54   $\times$   10 11   h/cm 2   ,   where 
a   safety   factor   of   3.0   for   simulation   was   included   in   the 
calculations.

V.   C ONCLUSIONS

The   radiation   tests   performed   have   qualified   the   key   com- 
ponents   of   the   DB   design   for   TID   and   NIEL   resistance   up 
to   the   required   radiation   levels.   Additionally,   SEE   tests   have 
allowed   to   eliminate   the   possibility   of   SEL   on   the   KU   FPGA 
used   in   the   current   revision   of   the   design   and   estimate   the 
impact of SEUs on the data taking for nominal runs. The team 
plans   conducting   radiation   tests   in   mixed   fields   and   additional 
SEU   tests   with   operational   firmware   to   further   evaluate   the 
impact   of   SEU   rates   on   board   performance.   Additionally, 
further radiation tests will be required during production phase 
to   account   for   variance   in   the   batches   during   the   production 
of   components.   The   tests   will   follow   the   ATLAS   Radiation 
Tolerant   Electronics   Proposed   Guideline.   The   project   will 
produce   930   DBs   of   which   896   will   be   installed   in   ATLAS 
for   the   HL-LHC.

IV.   D AUGHTERBOARD   RADIATION   QUALIFICATION 
STUDIES

A   diverse   set   of   radiation   campaigns   took   place   were 
commercial   development   boards,   dedicated   boards,   and   full 
four   assembled   DB   prototypes   that   were   irradiated   to   more 
than   108   Gy   and   14   $\times$   10 12   n/cm 2 .   These   tests   qualified 
the   Kintex   Ultrascale   FPGAs,   the   ProASIC   FPGAs,   the 
IS25LP128F,   IS25WP256   and   IS25LP512   flash   memories, 
DAN217T146CT Schottky diodes, 540BAA100M000BBG os- 
cillators,   and   INA333AIDRG   instrumentation   amplifiers.   The 
ProASIC FPGAs and IS25LP128F flash memories the showed 
some   after-effects   on   the   re-programmability   of   the   devices 
that   were   fully   recovered   in   all   devices   after   annealing   and 
have   no   impact   in   the   nominal   run   of   the   system.   Similarly, 
the   impact   of   ionizing   radiation   on   the   threshold   voltages   the 
FDFMA3N109   MOSFETs   is   acceptable   for   the   digital   use   of 
these   devices   on   the   design   a.   In   the   case   of   the   QUAD   os- 
cillators   LTC6902IMS   it   was   determined   that   the   inconsistent 
post   irradiation   behavior   over   the   sample   set   was   sufficient 
for   disqualifying   the   device.   Four   different   variants   of   SFP+ 
devices   from   four   different   companies   were   tested:   AVAGO, 
CORETEK,   FS-Europe   and   Direktronik.   The   AVAGO   AFBR- 
709SMZ the CORETEK CT-A000NPP-SB1L-D devices failed 
to   pass   the   irradiation   tests,   while   both   the   Direktronik   33- 
4248   and   the   FS-Europe   SFP-10GSR-85   were   qualified   with

R EFERENCES

[1]   ATLAS   Collaboration.   Technical   Design   Report   for   the   Phase-II   Up- 
grade of the ATLAS Tile Calorimeter. CERN-LHCC-2017-019, ATLAS- 
TDR-028,   2018 
[2]   E.   Valdes   Santurio.   Development   of   the   read-out   link   and   control 
board   for   the   ATLAS   Tile   Calorimeter   Upgrade,   PhD   dissertation, 
Department   of   Physics,   Stockholm   University,   Stockholm,   2019.   DiVA: 
diva2:1364842. 
[3]   GBTX 
Manual, 
https://espace.cern.ch/GBT-Project/GBTX/Manuals/ 
gbtxManual.pdf. 
[4]   C. 
Amelung, 
H. 
Chen, 
2020, 
ATLAS 
Radiation 
Tolerant 
Electronics   Proposed   Guideline   on   COTS   Lot   to   Lot   Variation 
https://indico.cern.ch/event/946481/contributions/3978191/attachments/ 
2087789/3507703/COTSL2L.20200806.v2.pdf 
[5]   K. 
ATLAS 
Collaboration, 
Radiation 
Environment 
of 
ATLAS. 
https://twiki.cern.ch/twiki/pub/Atlas/RadiationMapsGeant4/ 
WebRadMaps   Zoom   R2   public.html

\end{document}
